
\newcommand{\listing}[1]{Listing~\ref{lst:#1}}

\newcommand{\bs}{$\backslash$}

% erzeugt ein Listenelement mit fetter Überschrift 
\newcommand{\itemd}[2]{\item{\textbf{#1}}\\{#2}}

% Usage:
% Foo bar & \specialcell{Foo\\bar}    & Foo bar \\ % vertically centered
% Foo bar & \specialcell[t]{Foo\\bar} & Foo bar \\ % aligned with top rule
% Foo bar & \specialcell[b]{Foo\\bar} & Foo bar \\ % aligned with bottom rule
\newcommand{\specialcell}[2][c]{\begin{tabular}[#1]{@{}l@{}}#2\end{tabular}}

\newcommand{\theadcenter}[2]{\multicolumn{1}{#1}{\bfseries #2}}
\newcommand{\thead}[1]{\bfseries #1}

% see also http://jevopi.blogspot.de/2010/04/latex-lstinline-and-hyphenations.html
\newcommand{\codeJava}[1]{\lstinline[language=Java,breaklines=true,mathescape,literate={\-}{}{0\discretionary{-}{}{}}]§#1§}
\newcommand{\codeCpp}[1]{\lstinline[language=C++,breaklines=true,mathescape,literate={\-}{}{0\discretionary{-}{}{}}]§#1§}

\newcommand{\code}[1]{\lstinline$#1$}
\newcommand{\file}[1]{\texttt{#1}}

\newcommand{\TBD}{\textcolor{red}{TBD}\xspace}
\newcommand{\TBC}{\textcolor{red}{TBC}\xspace}
%\newcommand{\TBD}{{\hypersetup{linkcolor=red}\acs{tbd}}}
%\newcommand{\TBC}{{\hypersetup{linkcolor=red}\acs{tbd}}}

\newcommand{\applicableDocument}[1]{[\hyperlink{AD#1}{AD#1}]}
\newcommand{\referenceDocument}[1]{[\hyperlink{RD#1}{RD#1}]}

\newcommand{\imageAttention}{\includegraphics[width=14pt]{attention}}	% TODO
\newenvironment{important}[1]
	{\begin{bclogo}[logo=\imageAttention,noborder=true, couleur=gray!10, couleurBarre=gray!10]{#1\vspace{.3cm}}}
	{\end{bclogo}}
